\section{Outlook}\label{sec:outlook}

\subsection{What is settled}

Three things are now closed, and it is worth separating them from what is not.

\textbf{The generator and the transfer are identified.}
Theorem~\ref{thm:kernel} writes the form in closed shape and shows the whole
shift dependence is one hyperbolic factor; Proposition~\ref{prop:unimodular}
shows the transfer is an all-pass filter, exactly and unconditionally. Everything
in Sections~\ref{sec:allpass}--\ref{sec:region} follows from the second.

\textbf{Two of the program's own facts are explained rather than recorded.} The
density identity $b(\tau^2)+w_0=2\theta'(\tau)$ is the statement that the Weil
symbol is a group delay (Proposition~\ref{prop:groupdelay}), and the band edge at
$\gamma_1$ of the seventh selection rule is the statement that the cheapest way
to avoid every zero is to live below the lowest one (Section~\ref{sec:region}).

\textbf{The shift--length plane is settled.} Corollary~\ref{cor:graded} makes the
shifted family a graded criterion, one zero-free strip per shift; the contraction
region has no boundary under the Riemann hypothesis, so the critical path sought
by \cite{CriticalPath} does not exist; and Remark~\ref{rem:disproof} says what
the plane is for instead.

\subsection{What is open, ranked}

\begin{enumerate}
\item \textbf{The group-delay test on a specific open Wilson line.}
Condition~\ref{cond:bilinear} is one computation: take a protected open Wilson
line two-point function in the defect setting of \cite{OWL}, extract the phase of
its spectral kernel against the dilatation eigenvalue, and compare the derivative
with $2\theta'(\tau)$. A bounded delay kills the candidate; a delay growing like
$\log\tau$ with the wrong coefficient kills it too. This needs no arithmetic and
no positivity, and it is the first item because it is the first test in this
program that a gauge theory can fail.

\item \textbf{How does the zero-placement precision collapse?}
Proposition~\ref{prop:kappa} reduces the margin to $\varepsilon_L$, the accuracy
with which a function supported in $I_L$ can put its spectral zeros on the
Riemann zeros, and the recorded decay rules out the naive prolate estimate. This
is the companion investigation's request for a mechanism rather than a
measurement, now with a measured law attached.

\item \textbf{Read Suzuki \cite{SuzukiDeBranges}.} Three routes now point at it:
the resonance picture of Proposition~\ref{prop:resonance}, the positive-real
reading of Proposition~\ref{prop:cayley}, and the Cayley compactification of
Section~\ref{sec:conditions}. The companion investigation already ranks it first
for an independent reason.

\item \textbf{Problems~\ref{prob:trace} and \ref{prob:curvature}}, which
Section~\ref{sec:endpoint} shows are the two forms of one question.

\item \textbf{Does the vanishing theorem generalize?} Whether "BPS implies
vanishing expectation unless superconformal" extends beyond the class of
\cite{OWL}. If it does, it is a general obstruction of the same kind as this
program's own exclusions and belongs in the selection index; if it does not, the
exception is where to look.

\item Only after these, a specific theory. Section~\ref{sec:conditions} says what
it must satisfy, and Remark~\ref{rem:angle} says where to look first.
\end{enumerate}

\subsection{Scope}

Nothing above is proposed as a route to the Riemann hypothesis, and
Section~\ref{sec:region} explains structurally why no argument of this shape could
be one.

Written proofs: Theorem~\ref{thm:kernel}, Proposition~\ref{prop:unimodular},
Proposition~\ref{prop:groupdelay}, Proposition~\ref{prop:resonance},
Proposition~\ref{prop:flux}, Lemma~\ref{lem:multiplicative},
Theorem~\ref{thm:abelian}, Proposition~\ref{prop:hyperbolic},
Proposition~\ref{prop:monotone}, Proposition~\ref{prop:dichotomy},
Corollary~\ref{cor:graded}, Proposition~\ref{prop:cayley},
Proposition~\ref{prop:endpoint}, Proposition~\ref{prop:algebra} and
Proposition~\ref{prop:flat}.

Conditional on the Riemann hypothesis: Proposition~\ref{prop:flux}, and the
second case of Proposition~\ref{prop:dichotomy} uses the weaker hypothesis stated
there. Proposition~\ref{prop:resonance}'s passage from one zero to the sum is
stated rather than proved uniformly and is supported by
Appendix~\ref{app:checks}.

Proof sketch: Proposition~\ref{prop:trace}, resting on a smoothing estimate
quoted from \cite{Background} and on the Sobolev trace threshold.

Heuristic, and labelled where it occurs: the single-parameter relation
$\kappa_L\sqrt{m_L}\approx\sqrt{Q[Xe_L]}$ of Proposition~\ref{prop:kappa}, whose
support is three recorded data points; the detection scale $\omega\sim1/L$ of
Remark~\ref{rem:artifact}; and Remark~\ref{rem:principal}, which is a reading
resting on an identification the program declines to make.

Not claimed anywhere: a source, a positivity statement, a matched gauge theory,
an excluded gauge theory, or novelty for the idea that the zeros behave like
resonances.

\subsection*{Preparation}

This working manuscript was produced with a large language model, Anthropic's
Claude (model \texttt{claude-opus-5}), working directly in the investigation
repository: the derivations and proofs of
Sections~\ref{sec:generator}--\ref{sec:endpoint}, the conditions of
Section~\ref{sec:conditions}, the check program of Appendix~\ref{app:checks}, and
the drafting throughout. The proposal the paper analyses is Edward Baker's, as is
the open Wilson line construction of \cite{OWL} and the earlier shifted-zeta
material of \cite{CriticalPath,CumulativePath} on which Section~\ref{sec:region}
builds; Remark~\ref{rem:attribution} records the priority of the hyperbolic
identity. The research notes accompanying the investigation name the model used
for each and record the exploratory computations deliberately outside the
registered checks. Every statement is classified in the text as a written proof,
a proof sketch, a labelled numerical computation, a heuristic, or a reading.
